\documentclass[10pt,a4paper]{amsart}
\usepackage{amsmath,amssymb,amsthm,graphicx}
\usepackage{verbatim}
\newenvironment{reference}{\comment}{\endcomment}
%\newenvironment{reference}{}{}
\newenvironment{slogan}{\comment}{\endcomment}
\newenvironment{history}{\comment}{\endcomment}

% For commutative diagrams we use Xy-pic
\usepackage[all]{xy}

% We use 2cell for 2-commutative diagrams.
\xyoption{2cell}
% Theorem environments.
%
\theoremstyle{plain}
\newtheorem{theorem}[subsection]{Theorem}
\newtheorem{proposition}[subsection]{Proposition}
\newtheorem{lemma}[subsection]{Lemma}

\theoremstyle{definition}
\newtheorem{definition}[subsection]{Definition}
\newtheorem{example}[subsection]{Example}
\newtheorem{exercise}[subsection]{Exercise}
\newtheorem{situation}[subsection]{Situation}

\theoremstyle{remark}
\newtheorem{remark}[subsection]{Remark}
\newtheorem{remarks}[subsection]{Remarks}

\numberwithin{equation}{subsection}

% Macros
%
\def\lim{\mathop{\mathrm{lim}}\nolimits}
\def\colim{\mathop{\mathrm{colim}}\nolimits}
\def\Spec{\mathop{\mathrm{Spec}}}
\def\Hom{\mathop{\mathrm{Hom}}\nolimits}
\def\Ext{\mathop{\mathrm{Ext}}\nolimits}
\def\SheafHom{\mathop{\mathcal{H}\!\mathit{om}}\nolimits}
\def\SheafExt{\mathop{\mathcal{E}\!\mathit{xt}}\nolimits}
\def\Sch{\mathit{Sch}}
\def\Mor{\mathop{\mathrm{Mor}}\nolimits}
\def\Ob{\mathop{\mathrm{Ob}}\nolimits}
\def\Sh{\mathop{\mathit{Sh}}\nolimits}
\def\NL{\mathop{N\!L}\nolimits}
\def\CH{\mathop{\mathrm{CH}}\nolimits}
\def\proetale{{pro\text{-}\acute{e}tale}}
\def\etale{{\acute{e}tale}}
\def\QCoh{\mathit{QCoh}}
\def\Ker{\mathop{\mathrm{Ker}}}
\def\Im{\mathop{\mathrm{Im}}}
\def\Coker{\mathop{\mathrm{Coker}}}
\def\Coim{\mathop{\mathrm{Coim}}}

% Boxtimes
%
\DeclareMathSymbol{\boxtimes}{\mathbin}{AMSa}{"02}

%
% Macros for moduli stacks/spaces
%
\def\QCohstack{\mathcal{QC}\!\mathit{oh}}
\def\Cohstack{\mathcal{C}\!\mathit{oh}}
\def\Spacesstack{\mathcal{S}\!\mathit{paces}}
\def\Quotfunctor{\mathrm{Quot}}
\def\Hilbfunctor{\mathrm{Hilb}}
\def\Curvesstack{\mathcal{C}\!\mathit{urves}}
\def\Polarizedstack{\mathcal{P}\!\mathit{olarized}}
\def\Complexesstack{\mathcal{C}\!\mathit{omplexes}}
% \Pic is the operator that assigns to X its picard group, usage \Pic(X)
% \Picardstack_{X/B} denotes the Picard stack of X over B
% \Picardfunctor_{X/B} denotes the Picard functor of X over B
\def\Pic{\mathop{\mathrm{Pic}}\nolimits}
\def\Picardstack{\mathcal{P}\!\mathit{ic}}
\def\Picardfunctor{\mathrm{Pic}}
\def\Deformationcategory{\mathcal{D}\!\mathit{ef}}
\usepackage{cleveref}
\crefname{originalreference}{original reference}{original references}
\Crefname{originalreference}{Original reference}{Original references}
\makeatletter
\newlabel{categories-lemma-characterize-right-exact}{{0GMN}{source}{}{sourcepointer-1}{}}
\newlabel{categories-lemma-characterize-right-exact@cref}{{[originalreference][0][]0GMN}{[1][0][]source}{}{}{}}
\makeatother
\title{Right exact functors on finitely presented modules are equivalent to pairs (K, κ: A → End\_B(K))}
\author{The Stacks Project Authors}
\date{stacks-project:functors.tex:local-snapshot-20261012; GNU FDL 1.2 or later}
\begin{document}
\maketitle
\noindent\textit{Editorial note.} All delivered mathematical text is the author's own, copied verbatim from the pinned Stacks Project chapter file functors.tex; no statement or proof step was rewritten, reordered or removed.\par\medskip
\noindent\textit{Editorial note.} Span recomputation by label pairing (LANE\_BRIEF 10.2): in this chapter the author's proof environment follows the statement environment as a standalone block; the spec covers the complete environment pair (statement span, proof span) as recorded in the pool candidate.\par\medskip
\noindent\textit{Editorial note.} Occupation scan: combined claim\_key + statement\_anchor + generic-line scan over all registration files (active and withdrawn) found 0 hits for functors.tex (fu-occupied.json, 2026-10-14); consistent with the duty manager's six-check dispatch.\par\medskip
\ifdefined\hypertarget\hypertarget{sourcepointer-1}{}\fi
\noindent\textit{Source reference.} Original reference 0GMN: sources/stacks-chow/tags.txt, source lines 20763--20763. Cross-chapter standard prerequisite cited by the author inside the counted span; it is used as a pointer and its own proof is not part of this item.\par\medskip
\begin{lemma}
\label{lemma-functor-on-modules-fp}
Let $A$ be a ring. Let $\mathcal{B}$ be an additive category
with cokernels. There is an equivalence of categories between
\begin{enumerate}
\item the category of functors $F : \text{Mod}^{fp}_A \to \mathcal{B}$
which are right exact, and
\item the category of pairs $(K, \kappa)$ where $K \in \Ob(\mathcal{B})$
and $\kappa : A \to \text{End}_\mathcal{B}(K)$ is a ring homomorphism
\end{enumerate}
given by the rule sending $F$ to $F(A)$ with its natural $A$-action.
\end{lemma}

\begin{proof}
Let $(K, \kappa)$ be as in (2). We will construct a functor
$F : \text{Mod}^{fp}_A \to \mathcal{B}$ such that $F(A) = K$
endowed with the given $A$-action $\kappa$. Namely, given an
integer $n \geq 0$ let us set
$$
F(A^{\oplus n}) = K^{\oplus n}
$$
Given an $A$-linear map $\varphi : A^{\oplus m} \to A^{\oplus n}$
with matrix $(a_{ij}) \in \text{Mat}(n \times m, A)$ we define
$$
F(\varphi) :
F(A^{\oplus m}) = K^{\oplus m}
\longrightarrow
K^{\oplus n} = F(A^{\oplus n})
$$
to be the map with matrix $(\kappa(a_{ij}))$. This defines an additive
functor $F$ from the full subcategory of
$\text{Mod}^{fp}_A$ with objects $0$, $A$, $A^{\oplus 2}$, $\ldots$
to $\mathcal{B}$; we omit the verification.

\medskip\noindent
For each object $M$ of $\text{Mod}^{fp}_A$ choose a presentation
$$
A^{\oplus m_M} \xrightarrow{\varphi_M} A^{\oplus n_M} \to M \to 0
$$
of $M$ as an $A$-module. Let us use the trivial presentation
$0 \to A^{\oplus n} \xrightarrow{1} A^{\oplus n} \to 0$ if $M = A^{\oplus n}$
(this isn't necessary but simplifies the exposition).
For each morphism $f : M \to N$ of
$\text{Mod}^{fp}_A$ we can choose a commutative diagram
\begin{equation}
\label{equation-map}
\vcenter{
\xymatrix{
A^{\oplus m_M} \ar[r]_{\varphi_M} \ar[d]_{\psi_f} &
A^{\oplus n_M} \ar[r] \ar[d]_{\chi_f} &
M \ar[r] \ar[d]_f & 0 \\
A^{\oplus m_N} \ar[r]^{\varphi_N} &
A^{\oplus n_N} \ar[r] &
N \ar[r] & 0
}
}
\end{equation}
Having made these choices we can define: for an object
$M$ of $\text{Mod}^{fp}_A$ we set
$$
F(M) = \Coker(F(\varphi_M) : F(A^{\oplus m_M}) \to F(A^{\oplus n_M}))
$$
and for a morphism $f : M \to N$ of $\text{Mod}^{fp}_A$ we set
$$
F(f) = \text{the map }F(M) \to F(N)\text{ induced by }
F(\psi_f)\text{ and }F(\chi_f)\text{ on cokernels}
$$
Note that this rule extends the given functor $F$ on
the full subcategory consisting of the free modules $A^{\oplus n}$.
We still have to show that $F$ is a functor, that $F$ is additive,
and that $F$ is right exact.

\medskip\noindent
Let $f : M \to N$ be a morphism $\text{Mod}^{fp}_A$. We claim that the map
$F(f)$ defined above is independent of the choices of $\psi_f$ and $\chi_f$
in (\ref{equation-map}). Namely, say
$$
\xymatrix{
A^{\oplus m_M} \ar[r]_{\varphi_M} \ar[d]_\psi &
A^{\oplus n_M} \ar[r] \ar[d]_\chi &
M \ar[r] \ar[d]_f & 0 \\
A^{\oplus m_N} \ar[r]^{\varphi_N} &
A^{\oplus n_N} \ar[r] &
N \ar[r] & 0
}
$$
is also commutative. Denote $F(f)' : F(M) \to F(N)$ the map
induced by $F(\psi)$ and $F(\chi)$. Looking at the commutative
diagrams, by elementary commutative algebra there exists a map
$\omega : A^{\oplus n_M} \to A^{\oplus m_N}$ such that
$\chi = \chi_f + \varphi_N \circ \omega$. Applying $F$ we
find that $F(\chi) = F(\chi_f) + F(\varphi_N) \circ F(\omega)$.
As $F(N)$ is the cokernel of $F(\varphi_N)$ we find
that the map $F(A^{\oplus n_M}) \to F(M)$ equalizes $F(f)$ and $F(f)'$.
Since a cokernel is an epimorphism, we conclude that $F(f) = F(f)'$.

\medskip\noindent
Let us prove $F$ is a functor. First, observe that
$F(\text{id}_M) = \text{id}_{F(M)}$ because we may pick
the identities for $\psi_f$ and $\chi_f$ in the diagram above
in case $f = \text{id}_M$. Second, suppose we have
$f : M \to N$ and $g : L \to M$. Then we see that
$\psi = \psi_f \circ \psi_g$ and $\chi = \chi_f \circ \chi_g$
fit into (\ref{equation-map}) for $f \circ g$.
Hence these induce the correct map which exactly
says that $F(f) \circ F(g) = F(f \circ g)$.

\medskip\noindent
Let us prove that $F$ is additive. Namely, suppose we have
$f, g : M \to N$. Then we see that $\psi = \psi_f + \psi_g$ and
$\chi = \chi_f + \chi_g$ fit into (\ref{equation-map}) for $f + g$.
Hence these induce the correct map which exactly
says that $F(f) + F(g) = F(f + g)$.

\medskip\noindent
Finally, let us prove that $F$ is right exact. It suffices to show that $F$
commutes with coequalizers, see
Categories, Lemma \ref{categories-lemma-characterize-right-exact}.
For this, it suffices to prove that $F$ commutes with cokernels.
Let $K \to L \to M \to 0$ be an exact sequence of $A$-modules
with $K$, $L$, $M$ finitely presented. Since $F$ is an additive
functor, this certainly gives a complex
$$
F(K) \to F(L) \to F(M) \to 0
$$
and we have to show that the second arrow is the cokernel of the first
in $\mathcal{B}$. In any case, we obtain a map
$\Coker(F(K) \to F(L)) \to F(M)$.
By elementary commutative algebra there exists a commutative diagram
$$
\xymatrix{
A^{\oplus m_M} \ar[r]_{\varphi_M} \ar[d]_\psi &
A^{\oplus n_M} \ar[r] \ar[d]_\chi &
M \ar[r] \ar[d]_1 & 0 \\
K \ar[r] &
L \ar[r] &
M \ar[r] & 0
}
$$
Applying $F$ to this diagram and using the construction of $F(M)$ as
the cokernel of $F(\varphi_M)$ we find there exists a map
$F(M) \to \Coker(F(K) \to F(L))$ which is a right inverse
to the map $\Coker(F(K) \to F(L)) \to F(M)$. This first implies
that $F(L) \to F(M)$ is an epimorphism always. Next, the above shows
we have
$$
\Coker(F(K) \to F(L)) = F(M) \oplus E
$$
where the direct sum decomposition is compatible with both
$F(M) \to \Coker(F(K) \to F(L))$ and $\Coker(F(K) \to F(L)) \to F(M)$.
However, then the epimorphism $p : F(L) \to E$ becomes zero both
after composition with $F(K) \to F(L)$ and after composition
with $F(A^{n_M}) \to F(L)$. However, since $K \oplus A^{n_M} \to L$
is surjective (algebra argument omitted), we conclude that
$F(K \oplus A^{n_M}) \to F(L)$ is an epimorphism (by the above)
whence $E = 0$. This finishes the proof.
\end{proof}

\end{document}
